%% ma: L12-B03-0011
%% lop: 12
%% chuong: 1
%% bai: 3
%% dang: 12.3.9
%% loai: TLN
%% muc: VD
%% dap_an: "4052"
%% nguon: "Ảnh giáo viên gửi 11/10/2026 (Đề kiểm tra giữa HK1 Toán 12, THPT Việt Đức 2026-2027, mã đề 0103), Phần III Câu 6"
%% kien_thuc: "Bài 3 (tiệm cận xiên), giới hạn của hàm số tại vô cực (lớp 11)"
%% trang_thai: cho-duyet
%% ly_do: "Dạng mới đề xuất, chờ giáo viên duyệt."
\begin{ex}
Đường tiệm cận xiên của đồ thị hàm số $y=\sqrt[3]{8x^3+2x+1}$ có phương trình là $y=ax+b$ $(a\ne0)$. Biểu thức $T=2026a+2027b$ có giá trị bằng bao nhiêu?
\shortans{4052}
\loigiai{$a=\lim\limits_{x\to\pm\infty}\dfrac yx=\lim\limits_{x\to\pm\infty}\sqrt[3]{8+\dfrac2{x^2}+\dfrac1{x^3}}=2$.

$b=\lim\limits_{x\to\pm\infty}\left(y-2x\right)$. Đặt $u=\sqrt[3]{8x^3+2x+1}$, $v=2x$ thì $u-v=\dfrac{u^3-v^3}{u^2+uv+v^2}=\dfrac{2x+1}{u^2+uv+v^2}$, và $u^2+uv+v^2\sim12x^2$ nên $b=\lim\dfrac{2x+1}{12x^2}=0$.

Vậy tiệm cận xiên là $y=2x$, $T=2026\cdot2+2027\cdot0=4052$.}
\end{ex}
